\documentclass{article}
\usepackage{amsmath, amsthm, amssymb}
\usepackage[margin=1.2in]{geometry}

\newtheorem{proposition}{Proposition}
\newtheorem{theorem}{Theorem}
\newtheorem{corollary}{Corollary}
\theoremstyle{definition}
\newtheorem{definition}{Definition}
\newtheorem{remark}{Remark}

\title{Directional Homogeneity in Multi-Dimensional Competition}
\date{}

\begin{document}
\maketitle

\section*{Model}

$N$ firms compete on $K$ dimensions. Firm $i$ chooses an allocation
$x_i = (x_{i,1}, \ldots, x_{i,K}) \in \mathbb{R}_{\geq 0}^K$ subject to
$\sum_k x_{i,k} \leq B_i$. Its payoff is
\[
  U_i(x_i, x_{-i}) = \sum_{j \neq i}\sum_{k=1}^{K} w_k\,
    \sigma\!\bigl(f(x_{i,k}) - f(x_{j,k})\bigr),
\]
where $\sigma(z) = (1+e^{-z})^{-1}$, $f : \mathbb{R}_{\geq 0} \to \mathbb{R}$
is strictly increasing and strictly concave, and $w \in \Delta^{K-1}$ are
evaluator weights.  At an interior Nash equilibrium, the KKT condition for
firm $i$ is
\begin{equation}\label{eq:kkt}
  w_k\,f'(x_{i,k})\,S_{i,k} = \lambda_i \quad \forall k,
\end{equation}
where $\lambda_i > 0$ is the budget multiplier and
\[
  S_{i,k} = \sum_{j \neq i} \sigma'\!\bigl(f(x_{i,k}) - f(x_{j,k})\bigr)
\]
is the competitive pressure firm $i$ faces on dimension $k$, with
$\sigma'(z) = \sigma(z)(1-\sigma(z)) \in (0,\tfrac{1}{4}]$.

\begin{definition}
The Nash equilibrium is \emph{directionally homogeneous} if there exists a
direction $p \in \Delta^{K-1}$ such that $x_{i,k}^* = p_k B_i$ for all $i, k$.
\end{definition}

\section{Exact Homogeneity Results}

\begin{proposition}[Uniform budgets]\label{prop:B}
If $B_i = B$ for all $i$, the unique Nash equilibrium satisfies
$x_{i,k}^* = x_{j,k}^*$ for all $i, j$, hence is directionally homogeneous.
\end{proposition}

\begin{proof}
The game is symmetric under permutation of firms: $U_i$ depends on firm $i$
only through its budget, which is the same for all $i$.  Since $f$ is strictly
concave, $U_i$ is strictly concave in $x_i$, so the Nash equilibrium is unique
(the best-response map is a strict contraction on the product space of budget
simplices).  Any permutation of a Nash equilibrium is again a Nash equilibrium;
by uniqueness, the equilibrium must be invariant under all firm permutations, so
$x_{i,k}^* = x_k^*$ for all $i$.  At this symmetric profile the competitive
pressures are $S_{i,k} = (N-1)\sigma'(0) = (N-1)/4$ for all $i, k$, and
\eqref{eq:kkt} becomes the standard water-filling equation $w_k f'(x_k^*) = 4\lambda/(N-1)$.
The common direction $p_k = x_k^*/B$ satisfies $\sum_k p_k = 1$, so the
equilibrium is directionally homogeneous.
\end{proof}

\begin{proposition}[Uniform weights]\label{prop:w}
If $w_k = 1/K$ for all $k$, then $x_{i,k}^* = B_i/K$ for all $i, k$ is a Nash
equilibrium, which is directionally homogeneous with direction $p_k = 1/K$.
\end{proposition}

\begin{proof}
We verify the KKT conditions \eqref{eq:kkt} at the proposed profile.  Since
$x_{i,k} = B_i/K$ is the same for all $k$ within each firm, the output gap
between any two firms is $k$-independent:
\[
  f(x_{i,k}) - f(x_{j,k}) = f(B_i/K) - f(B_j/K) =: \delta_{ij}.
\]
Hence $S_{i,k} = \sum_{j \neq i} \sigma'(\delta_{ij}) =: S_i$ is also $k$-independent.
The KKT condition becomes $\tfrac{1}{K} f'(B_i/K)\, S_i = \lambda_i$ for all $k$,
which is satisfied with $\lambda_i = f'(B_i/K) S_i / K > 0$.  The budget binds:
$\sum_k B_i/K = B_i$.  Hence the profile is a Nash equilibrium, and
$x_{i,k}^*/B_i = 1/K$ is the same for all firms.
\end{proof}

\section{Characterization via $\Omega_f$}

Define the \emph{scale-adjusted marginal product} $\Omega_f(x) = x\,f'(x)$.

\begin{theorem}[$\Omega_f$ characterization]\label{thm:log}
The allocation $x_{i,k}^* = w_k B_i$ is a Nash equilibrium for all budget
vectors $B$ and weight vectors $w$ if and only if $f(x) = a\log x + b$ for
some constants $a > 0$, $b$---equivalently, $\Omega_f(x)$ is constant.
\end{theorem}

\begin{proof}
\textbf{($\Leftarrow$)} Suppose $f(x) = a\log x + b$, so $f'(x) = a/x$ and
$\Omega_f \equiv a$.  At $x_{i,k}^* = w_k B_i$, the output gap is
\[
  f(w_k B_i) - f(w_k B_j) = a\log\frac{w_k B_i}{w_k B_j}
    = a\log\frac{B_i}{B_j},
\]
which is $k$-independent.  Hence $S_{i,k} = \sum_{j \neq i}\sigma'(a\log(B_i/B_j))$
is $k$-independent.  The KKT condition becomes
\[
  w_k \cdot \frac{a}{w_k B_i} \cdot S_i = \frac{a S_i}{B_i} = \lambda_i \quad\forall k.
\]
The budget binds: $\sum_k w_k B_i = B_i$.  Hence the profile is a Nash equilibrium
for any $B$ and $w$.

\textbf{($\Rightarrow$)} Suppose $x_{i,k}^* = p_k B_i$ is a Nash equilibrium
for all $B$ and $w$.  Fix any two firms with $B_1 \neq B_2$ and consider the
KKT conditions \eqref{eq:kkt} for these two firms:
\[
  w_k\,f'(p_k B_1)\,S_{1,k} = \lambda_1, \qquad
  w_k\,f'(p_k B_2)\,S_{2,k} = \lambda_2.
\]
For $N = 2$, $S_{1,k} = \sigma'(\Delta_k)$ and $S_{2,k} = \sigma'(-\Delta_k)
= \sigma'(\Delta_k)$ (since $\sigma'$ is even), where $\Delta_k = f(p_k B_1)
- f(p_k B_2)$.  Dividing the two KKT conditions:
\begin{equation}\label{eq:ratio}
  \frac{f'(p_k B_1)}{f'(p_k B_2)} = \frac{\lambda_1}{\lambda_2} \quad\forall k.
\end{equation}
This ratio is $k$-independent.  Since $p_k > 0$ can take any positive value
as $w$ varies over the interior of $\Delta^{K-1}$, equation \eqref{eq:ratio}
says $f'(t B_1)/f'(t B_2)$ is constant in $t > 0$ for all $B_1, B_2 > 0$.
Differentiating in $t$:
\[
  B_1\,\frac{f''(t B_1)}{f'(t B_1)} = B_2\,\frac{f''(t B_2)}{f'(t B_2)}
  \quad\forall t, B_1, B_2.
\]
Setting $x = t B_1$ and $y = t B_2$ (any $x, y > 0$ are achieved by varying
$t$ and the ratio $B_1/B_2$), this gives $x f''(x)/f'(x) = $ const, i.e., $\eta_f$
is constant.  Thus $f'(x) = A x^\beta$ for some $A > 0$, $\beta$.

For $\beta \neq -1$: $f(x) \propto x^{\beta+1}$ and the output gap is
$\Delta_k = \tfrac{A}{\beta+1} p_k^{\beta+1}(B_1^{\beta+1} - B_2^{\beta+1})$,
which depends on $k$ through $p_k$.  Then $S_{1,k} = \sigma'(\Delta_k)$ is
$k$-dependent, so the KKT condition $w_k f'(p_k B_1) S_{1,k} = \lambda_1$ fails
for non-uniform $w$ (since $w_k f'(p_k B_1) = w_k A(p_k B_1)^\beta$ is $k$-independent
only for the specific direction $p_k \propto w_k^{-1/\beta}$, but then $S_{1,k}$
still varies unless $\Delta_k$ is constant).  Hence $\beta = -1$, i.e., $f(x) = a\log x + b$
and $\Omega_f(x) = a$.
\end{proof}

\begin{remark}[Shifted log]
The family $f(x) = a\log(x + c) + b$ with $c \geq 0$ also gives exact
directional homogeneity under a slightly modified notion: the Nash is
$x_{i,k}^* = w_k(B_i + cK) - c$, i.e., the shifted allocations
$\tilde x_{i,k} = x_{i,k} + c = w_k(B_i + cK)$ are proportional.  The
same argument shows this is the unique affinely-proportional family.
\end{remark}

\begin{corollary}[Comparative statics]\label{cor:comp}
For two production functions $h_1$ and $h_2$, the Nash under $h_2$ is
more directionally homogeneous than under $h_1$ if $|\eta_{h_2}(x) - 1|
< |\eta_{h_1}(x) - 1|$ over the relevant investment range, i.e.,
$\Omega_{h_2}$ is more constant than $\Omega_{h_1}$.  In particular, for
power laws $f(x) = x^\alpha$ with $\alpha \in (0, 1)$:
$\eta_f = 1 - \alpha$ is constant with $|\eta_f - 1| = \alpha$, so higher
$\alpha$ produces more heterogeneous Nash equilibria, and the Nash converges
to the proportional allocation $x_{i,k}^* = w_k B_i$ as $\alpha \to 0$.
\end{corollary}

\begin{proof}
By Theorem~\ref{thm:log}, the deviation from directional homogeneity is
captured by how far $\Omega_f$ is from being constant.  For a proportional
perturbation around an approximately-homogeneous Nash with direction $p$,
the output gap between firms $i$ and $j$ is $f(p_k B_i) - f(p_k B_j)$.
Differentiating with respect to $\log p_k$:
\[
  \frac{d}{d\log p_k}\bigl[f(p_k B_i) - f(p_k B_j)\bigr]
  = \Omega_f(p_k B_i) - \Omega_f(p_k B_j).
\]
The gap is $k$-independent if and only if this derivative vanishes, which holds
iff $\Omega_f$ is constant (Theorem~\ref{thm:log}).  The magnitude of $k$-variation
is therefore controlled by $\mathrm{osc}(\Omega_f)$ over the relevant investment range.
For $f = x^\alpha$, $\Omega_f(x) = \alpha x^\alpha$ grows monotonically in $x$ with
$|1 - \eta_f| = \alpha$, while for $\alpha \to 0$ the function $x^\alpha \to 1$
approaches the log family in the sense that $\Omega_f \to $ const.
\end{proof}

\end{document}
